\begin{ExerciseBlock}
\ExerciseHeading{04}{Drehscheiben}{rigid_bodies_coupled_rotating_disks}
In der Abbildung rotiert eine Scheibe $A$ mit Winkelgeschwindigkeit $\omega_0 = 1 \,s^{-1}$ um die Stange. Es wird eine weitere, anfangs nicht rotierende, Scheibe $B$ des gleich Materials aufgesetzt.

Daten: $h^{(A)} = h^{(B)}$; $R_i^{(A)} = R_i^{(B)} = 5\,$cm; $R_a^{(A)} = 30\,$cm; $R_a^{(B)} = 20\,$cm;

\begin{enumerate}
	\item Bestimmen Sie die daraus resultierende Winkelgeschwindigkeit $\omega_1$.
	\item Wieviel Prozent $P$ der Rotationsenergie wird hierbei in  Wärmenergie $\Delta Q$ umgesetzt.

\end{enumerate}

\begin{center}
	\CachedTikzFigure{content/resources/chapters/rigid_bodies/tikz/coupled_rotating_disks.tex}
\end{center}
\end{ExerciseBlock}
\UseExerciseSolution{rigid_bodies_coupled_rotating_disks}
